\documentclass[12pt]{article}

\usepackage[T1]{fontenc}
\usepackage[utf8]{inputenc}
\usepackage[brazil]{babel}
\usepackage{lmodern}
\usepackage{microtype}
\usepackage[a4paper,margin=2.5cm]{geometry}
\usepackage{setspace}
\usepackage{indentfirst}

\setstretch{1.15}

\title{O FIM DO LAISSEZ-FAIRE \thanks{Texto Traduzido livremente por Luiz Mario Andrade com o auxílio da ferramenta Chat GPT 5.4}}
\author{John Maynard Keynes}
\date{Julho de 1926}

\begin{document}

\maketitle

Este ensaio, que foi publicado como um panfleto pela Hogarth Press em julho de 1926, baseou-se na Palestra Sidney Ball proferida por Keynes em Oxford, em novembro de 1924, e em uma palestra proferida por ele na Universidade de Berlim, em junho de 1926. Os capítulos iv e v foram utilizados em \textit{Ensaios de Persuasão}.

\begin{center}
\textbf{§ I}
\end{center}

A disposição em relação aos assuntos públicos, que convenientemente resumimos como individualismo e \textit{laissez-faire}, extraiu seu sustento de muitos riachos diferentes de pensamento e fontes de sentimento. Por mais de cem anos, nossos filósofos nos governaram porque, por um milagre, quase todos concordaram ou pareceram concordar com esta única coisa. Ainda não dançamos conforme uma nova melodia. Mas uma mudança está no ar. Ouvimos apenas indistintamente o que antes eram as vozes mais claras e distinguíveis que já instruíram a humanidade política. A orquestra de instrumentos diversos, o coro de sons articulados, está finalmente recuando para a distância.

No final do século XVII, o direito divino dos monarcas deu lugar à liberdade natural e ao contrato, e o direito divino da igreja ao princípio da tolerância e à visão de que uma igreja é "uma sociedade voluntária de homens", reunindo-se de uma forma que é "absolutamente livre e espontânea".\footnote{Locke, \textit{A Letter Concerning Toleration}.} Cinquenta anos depois, a origem divina e a voz absoluta do dever deram lugar aos cálculos de utilidade. Nas mãos de Locke e Hume, essas doutrinas fundaram o Individualismo. O contrato presumia direitos no indivíduo; a nova ética, sendo nada mais do que um estudo científico das consequências do amor-próprio racional, colocou o indivíduo no centro. "O único trabalho que a Virtude exige", disse Hume, "é o de um Cálculo justo e uma preferência constante pela maior Felicidade".\footnote{\textit{An Enquiry Concerning the Principles of Morals}, seção lx.} Essas ideias harmonizavam-se com as noções práticas de conservadores e de advogados. Elas forneceram um fundamento intelectual satisfatório para os direitos de propriedade e para a liberdade do indivíduo na posse de fazer o que quisesse consigo mesmo e com o que era seu. Essa foi uma das contribuições do século XVIII para o ar que ainda respiramos.

O propósito de promover o indivíduo era depor o monarca e a igreja; o efeito — através do novo significado ético atribuído ao contrato — foi fortalecer a propriedade e a prescrição. Mas não demorou muito para que as reivindicações da sociedade se levantassem novamente contra o indivíduo. Paley e Bentham aceitaram o hedonismo utilitarista\footnote{"Eu omito", diz o Arquidiácono Paley, "muita declamação usual sobre a dignidade e a capacidade de nossa natureza, a superioridade da alma sobre o corpo, da parte racional sobre a animal de nossa constituição; sobre o valor, refinamento e delicadeza de algumas satisfações, e a mesquinhez, grosseria e sensualidade de outras: porque sustento que os prazeres não diferem em nada além da continuidade e intensidade" (\textit{Principles of Moral and Political Philosophy}, Livro I, cap. 6).} das mãos de Hume e seus predecessores, mas o ampliaram em utilidade social. Rousseau pegou o Contrato Social de Locke e extraiu dele a Vontade Geral. Em cada caso, a transição foi feita em virtude da nova ênfase dada à igualdade. "Locke aplica seu Contrato Social para modificar a igualdade natural da humanidade, na medida em que essa frase implica igualdade de propriedade ou mesmo de privilégio, em consideração à segurança geral. Na versão de Rousseau, a igualdade não é apenas o ponto de partida, mas o objetivo".\footnote{Leslie Stephen, \textit{English Thought in the Eighteenth Century}, II, 192.}

Paley e Bentham chegaram ao mesmo destino, mas por caminhos diferentes. Paley evitou uma conclusão egoísta para seu hedonismo através de um \textit{Deus ex machina}. "Virtude", diz ele, "é fazer o bem à humanidade, em obediência à vontade de Deus, e por causa da felicidade eterna" — desta forma trazendo de volta o \textit{Eu} e os \textit{outros} a uma paridade. Bentham alcançou o mesmo resultado pela razão pura. Não há fundamento racional, argumentou ele, para preferir a felicidade de um indivíduo, mesmo a si próprio, à de qualquer outro. Portanto, a maior felicidade do maior número é o único objeto racional de conduta — adotando a utilidade de Hume, mas esquecendo o corolário cínico daquele homem sábio: "'Não é contrário à razão preferir a destruição do mundo inteiro a um arranhão no meu dedo. Não é contrário à razão eu escolher minha ruína total para evitar o menor desconforto de um índio, ou de uma pessoa totalmente desconhecida para mim... A razão é e deve ser apenas a escrava das paixões, e nunca pode pretender qualquer outro cargo senão o de servi-las e obedecê-las."

Rousseau derivou a igualdade do estado de natureza, Paley da vontade de Deus, Bentham de uma lei matemática de indiferença. Igualdade e altruísmo haviam assim entrado na filosofia política e, de Rousseau e Bentham em conjunto, surgiram tanto a democracia quanto o socialismo utilitarista.

Esta é a segunda corrente — surgida de controvérsias mortas há muito tempo e levada em seu caminho por sofismas há muito explorados — que ainda permeia nossa atmosfera de pensamento. Mas ela não expulsou a primeira corrente. Misturou-se a ela. O início do século XIX realizou a união milagrosa. Harmonizou o individualismo conservador de Locke, Hume, Johnson e Burke com o socialismo e o igualitarismo democrático de Rousseau, Paley, Bentham e Godwin.\footnote{Godwin levou o \textit{laissez-faire} tão longe que pensou que todo governo era um mal, no que Bentham quase concordou com ele. A doutrina da igualdade torna-se com ele uma de individualismo extremo, beirando a anarquia. "O exercício universal do julgamento privado", diz ele, "é uma doutrina tão indizivelmente bela que o verdadeiro político certamente sentirá infinita relutância em admitir a ideia de interferir nela" (ver Leslie Stephen, \textit{op. cit.} II, 477).}

No entanto, aquela era teria tido dificuldade em alcançar essa harmonia de opostos se não fosse pelos economistas, que ganharam destaque no momento exato. A ideia de uma harmonia divina entre a vantagem privada e o bem público já é aparente em Paley. Mas foram os economistas que deram à noção uma boa base científica. Suponha que, pelo funcionamento das leis naturais, os indivíduos que buscam seus próprios interesses com esclarecimento em condições de liberdade sempre tendam a promover o interesse geral ao mesmo tempo! Nossas dificuldades filosóficas estão resolvidas — pelo menos para o homem prático, que pode então concentrar seus esforços em assegurar as condições necessárias de liberdade. À doutrina filosófica de que o governo não tem o direito de interferir, e à divina de que não tem necessidade de interferir, acrescenta-se uma prova científica de que sua interferência é desnecessária e inoportuna. Esta é a terceira corrente de pensamento, apenas descobrível em Adam Smith, que estava pronto, no geral, a permitir que o bem público repousasse no "esforço natural de cada indivíduo para melhorar sua própria condição", mas não totalmente e conscientemente desenvolvida até que o século XIX começasse. O princípio do \textit{laissez-faire} havia chegado para harmonizar o individualismo e o socialismo, e para tornar um só o egoísmo de Hume com o maior bem do maior número. O filósofo político poderia se aposentar em favor do homem de negócios — pois este último poderia alcançar o \textit{summum bonum} do filósofo apenas buscando seu próprio lucro privado.

No entanto, alguns outros ingredientes eram necessários para completar o pudim. Primeiro, a corrupção e a incompetência do governo do século XVIII, muitos legados dos quais sobreviveram no século XIX. O individualismo dos filósofos políticos apontava para o \textit{laissez-faire}. A harmonia divina ou científica (conforme o caso) entre o interesse privado e a vantagem pública apontava para o \textit{laissez-faire}. Mas, acima de tudo, a inépcia dos administradores públicos prejudicou fortemente o homem prático em favor do \textit{laissez-faire} — um sentimento que de forma alguma desapareceu. Quase tudo o que o Estado fez no século XVIII além de suas funções mínimas foi, ou pareceu, prejudicial ou mal-sucedido.

Por outro lado, o progresso material entre 1750 e 1850 veio da iniciativa individual e não deveu quase nada à influência diretiva da sociedade organizada como um todo. Assim, a experiência prática reforçou os raciocínios \textit{a priori}. Os filósofos e os economistas nos disseram que, por diversas razões profundas, a livre iniciativa privada promoveria o maior bem de todos. O que poderia servir melhor ao homem de negócios? E poderia um observador prático, olhando ao seu redor, negar que as bênçãos de melhoria que distinguiram a era em que viveu eram rastreáveis às atividades de indivíduos "em ascensão"?

Assim, o terreno era fértil para uma doutrina de que, fosse por motivos divinos, naturais ou científicos, a ação estatal deveria ser estreitamente confinada e a vida econômica deixada, tanto quanto possível desregulada, à habilidade e ao bom senso de cidadãos individuais motivados pelo admirável motivo de tentar progredir no mundo.

Na época em que a influência de Paley e seus semelhantes estava minguando, as inovações de Darwin estavam abalando os fundamentos da crença. Nada poderia parecer mais oposto do que a velha doutrina e a nova — a doutrina que via o mundo como a obra do relojoeiro divino e a doutrina que parecia extrair todas as coisas do Acaso, do Caos e do Tempo Antigo. Mas neste ponto as novas ideias reforçaram as antigas. Os economistas ensinavam que a riqueza, o comércio e as máquinas eram filhos da livre concorrência — que a livre concorrência construiu Londres. Mas os darwinistas podiam ir além disso — a livre concorrência havia construído o homem. O olho humano não era mais a demonstração do design, contatando milagrosamente todas as coisas para o melhor; era a conquista suprema do acaso, operando sob condições de livre concorrência e \textit{laissez-faire}. O princípio da sobrevivência do mais apto poderia ser considerado como uma vasta generalização da economia ricardiana. Interferências socialistas tornaram-se, à luz desta síntese mais grandiosa, não apenas inoportunas, mas ímpias, como calculadas para retardar o movimento progressivo do poderoso processo pelo qual nós mesmos surgimos como Afrodite da lama primitiva do oceano.

Portanto, traço a unidade peculiar da filosofia política cotidiana do século XIX ao sucesso com que harmonizou escolas diversificadas e guerreiras e uniu todas as coisas boas a um único fim. Hume e Paley, Burke e Rousseau, Godwin e Malthus, Cobbett e Huskisson, Bentham e Coleridge, Darwin e o Bispo de Oxford, todos estavam, descobriu-se, pregando praticamente a mesma coisa — individualismo e \textit{laissez-faire}. Esta era a Igreja da Inglaterra e aqueles seus apóstolos, enquanto a companhia dos economistas estava lá para provar que o menor desvio para a impiedade envolvia a ruína financeira.

Essas razões e esta atmosfera são as explicações, quer saibamos ou não — e a maioria de nós nestes dias degenerados é amplamente ignorante no assunto — de por que sentimos um preconceito tão forte em favor do \textit{laissez-faire}, e por que a ação estatal para regular o valor do dinheiro, ou o curso do investimento, ou a população, provoca suspeitas tão apaixonadas em muitos peitos íntegros. Não lemos esses autores; consideraríamos seus argumentos ultrajantes se caíssem em nossas mãos. No entanto, não pensaríamos como pensamos, imagino, se Hobbes, Locke, Hume, Rousseau, Paley, Adam Smith, Bentham e Miss Martineau não tivessem pensado e escrito como pensaram. Um estudo da história da opinião é um preliminar necessário para a emancipação da mente. Não sei o que torna um homem mais conservador — não saber nada além do presente, ou nada além do passado.

\begin{center}
\textbf{II}
\end{center}

Eu disse que foram os economistas que forneceram o pretexto científico pelo qual o homem prático pôde resolver a contradição entre egoísmo e socialismo que emergiu da filosofagem do século XVIII e da decadência da religião revelada. Mas, tendo dito isso por brevidade, apresso-me a qualificá-lo. Isso é o que os economistas \textit{supostamente} disseram. Nenhuma doutrina desse tipo é realmente encontrada nos escritos das maiores autoridades. É o que os popularizadores e vulgarizadores disseram. É o que os Utilitaristas, que admitiam o egoísmo de Hume e o igualitarismo de Bentham ao mesmo tempo, foram levados a acreditar, se quisessem realizar uma síntese.\footnote{Pode-se simpatizar com a visão de Coleridge, conforme resumida por Leslie Stephen, de que "os Utilitaristas destruíram todo elemento de coesão, fizeram da Sociedade uma luta de interesses egoístas e atingiram as próprias raízes de toda ordem, patriotismo, poesia e religião".} A linguagem dos economistas prestou-se à interpretação do \textit{laissez-faire}. Mas a popularidade da doutrina deve ser atribuída aos filósofos políticos da época, a quem ela convinha, em vez de aos economistas políticos.

A máxima \textit{laissez-nous faire} é tradicionalmente atribuída ao comerciante Legendre dirigindo-se a Colbert em algum momento no final do século XVII.\footnote{"O que devemos fazer para ajudá-lo?", perguntou Colbert. "\textit{Nous laisser faire}", respondeu Legendre.} Mas não há dúvida de que o primeiro escritor a usar a frase, e a usá-la em associação clara com a doutrina, é o Marquês d'Argenson por volta de 1751.\footnote{Para a história da frase, ver Oncken, "\textit{Die Maxime Laissez faire et laissez passer}", de quem a maioria das citações a seguir é retirada. As reivindicações do Marquês d'Argenson foram negligenciadas até que Oncken as apresentou, em parte porque as passagens relevantes publicadas durante sua vida eram anônimas (\textit{Journal Économique}, 1751), e em parte porque suas obras não foram publicadas na íntegra (embora provavelmente circulassem privadamente durante sua vida) até 1858 (\textit{Mémoires et Journal inédit du Marquis d'Argenson}).} O Marquês foi o primeiro homem a se apaixonar pelas vantagens econômicas de os governos deixarem o comércio em paz. Para governar melhor, disse ele, deve-se governar menos.\footnote{"\textit{Pour gouverner mieux, il faudrait gouverner moins.}"} A verdadeira causa do declínio de nossas manufaturas, declarou ele, é a proteção que lhes demos.\footnote{"\textit{On ne peut dire autant de nos fabriques: la vraie cause de leur déclin, c'est la protection outrée qu'on leur accorde.}"} "\textit{Laissez faire, telle devrait être la devise de toute puissance publique, depuis que le monde est civilisé.}" "\textit{Détestable principe que celui de ne vouloir grandeur que par l'abaissement de nos voisins! Il n'y a que la méchanceté et la malignité du cœur de satisfaites dans ce principe, et l'intérêt y est opposé. Laissez faire, morbleu! Laissez faire!!}"

Aqui temos a doutrina econômica do \textit{laissez-faire}, com sua expressão mais fervorosa no livre comércio, totalmente vestida. As frases e a ideia devem ter passado correntes em Paris a partir dessa época. Mas demoraram a se estabelecer na literatura; e a tradição que as associa aos fisiocratas, e particularmente a de Gournay e Quesnay, encontra pouco apoio nos escritos desta escola, embora fossem, é claro, proponentes da harmonia essencial dos interesses sociais e individuais. A frase \textit{laissez-faire} não é encontrada nas obras de Adam Smith, de Ricardo ou de Malthus. Nem mesmo a ideia está presente em uma forma dogmática em qualquer um desses autores. Adam Smith, é claro, era um Livre-Comerciante e um oponente de muitas restrições do século XVIII ao comércio. Mas sua atitude em relação aos Atos de Navegação e às leis de usura mostra que ele não era dogmático. Até mesmo sua famosa passagem sobre a "mão invisível" reflete a filosofia que associamos a Paley, em vez do dogma econômico do \textit{laissez-faire}. Como Sidgwick e Cliff Leslie apontaram, a defesa de Adam Smith do "óbvio e simples sistema de liberdade natural" deriva de sua visão teísta e otimista da ordem do mundo, conforme estabelecido em sua \textit{Teoria dos Sentimentos Morais}, e não de qualquer proposição de economia política propriamente dita.\footnote{Sidgwick, \textit{Principles of Political Economy}, p. 20.} A frase \textit{laissez-faire} foi, creio eu, trazida pela primeira vez ao uso popular na Inglaterra por uma passagem bem conhecida do Dr. Franklin.\footnote{Bentham usa a expressão "\textit{laissez-nous faire}" (\textit{Works}, p. 440).} Não é, de fato, até chegarmos às obras posteriores de Bentham — que não era de forma alguma um economista — que descobrimos a regra do \textit{laissez-faire}, na forma em que nossos avós a conheciam, adotada a serviço da filosofia Utilitarista. Por exemplo, em \textit{A Manual of Political Economy},\footnote{Escrito em 1793, um capítulo publicado na \textit{Bibliothèque Britannique} em 1798, e o todo impresso pela primeira vez na edição de Bowring de suas \textit{Works} (1843).} ele escreve: "A regra geral é que nada deve ser feito ou tentado pelo governo; o lema ou palavra de ordem do governo, nestas ocasiões, deve ser — \textit{Fique quieto}... O pedido que a agricultura, a manufatura e o comércio apresentam aos governos é tão modesto e razoável quanto o que Diógenes fez a Alexandre: Saia da frente do meu sol".

A partir desse momento, foi a campanha política pelo livre comércio, a influência da chamada Escola de Manchester e dos Utilitaristas Benthamitas, os pronunciamentos de autoridades econômicas secundárias e as histórias educativas de Miss Martineau e da Sra. Marcet que fixaram o \textit{laissez-faire} na mente popular como a conclusão prática da economia política ortodoxa — com esta grande diferença: que a visão malthusiana da população tendo sido aceita entretanto por essa mesma escola de pensamento, o \textit{laissez-faire} otimista da última metade do século XVIII dá lugar ao \textit{laissez-faire} pessimista da primeira metade do século XIX.\footnote{Cf. Sidgwick, \textit{op. cit.} p. 22: "Mesmo aqueles economistas que aderiram no essencial às limitações de Adam Smith à esfera do governo, aplicaram essas limitações tristemente em vez de triunfantemente; não como admiradores da ordem social resultante da 'liberdade natural', mas como convencidos de que ela é pelo menos preferível a qualquer ordem artificial que o governo pudesse substituir por ela."}

Em \textit{Conversations on Political Economy} (1817) da Sra. Marcet, Caroline resiste o quanto pode em favor do controle dos gastos dos ricos. Mas na página 418 ela tem que admitir a derrota:

\begin{quote}
CAROLINE. Quanto mais aprendo sobre este assunto, mais me sinto convencida de que os interesses das nações, assim como os dos indivíduos, longe de se oporem uns aos outros, estão em perfeita harmonia.

SRA. B. Visões liberais e ampliadas sempre levarão a conclusões semelhantes, e nos ensinarão a cultivar sentimentos de benevolência universal uns para com os outros; daí a superioridade da ciência sobre o mero conhecimento prático.
\end{quote}

Por volta de 1850, as \textit{Easy Lessons for the Use of Young People}, do Arcebispo Whately, que a Sociedade para a Promoção do Conhecimento Cristão estava distribuindo por atacado, não admitem nem mesmo as dúvidas que a Sra. B. permitia que Caroline ocasionalmente alimentasse. "É provável que se faça mais mal do que bem", conclui o livrinho, "por quase qualquer interferência do Governo nas transações monetárias dos homens, seja arrendamento, ou compra e venda de qualquer tipo". A verdadeira liberdade é "que cada homem deve ser deixado livre para dispor de sua própria propriedade, seu próprio tempo, força e habilidade, da maneira que ele próprio achar adequada, desde que não prejudique seus vizinhos".

Em suma, o dogma havia se apoderado da máquina educacional; tornara-se uma máxima de caderno escolar. A filosofia política, que os séculos XVII e XVIII haviam forjado para derrubar reis e prelados, tornara-se leite para bebês e entrara literalmente no berçário.

Finalmente, nas obras de Bastiat chegamos à expressão mais extravagante e rapsódica da religião do economista político. Em suas \textit{Harmonies Économiques},

\begin{quote}
Eu me proponho [diz ele] a demonstrar a Harmonia daquelas leis da Providência que governam a sociedade humana. O que torna essas leis harmoniosas e não discordantes é que todos os princípios, todos os motivos, todas as fontes de ação, todos os interesses, cooperam para um grande resultado final... E esse resultado é a aproximação indefinida de todas as classes em direção a um nível, que está sempre subindo; em outras palavras, a equalização dos indivíduos na melhora geral.
\end{quote}

E quando, como outros sacerdotes, ele redige seu \textit{Credo}, ele reza o seguinte:

\begin{quote}
Creio que Aquele que organizou o universo material não retirou Seu olhar dos arranjos do mundo social. Creio que Ele combinou e fez mover em harmonia agentes livres, bem como moléculas inertes... Creio que a tendência social invencível é uma aproximação constante dos homens em direção a um nível moral, intelectual e físico comum, com, ao mesmo tempo, uma elevação progressiva e indefinida desse nível. Creio que tudo o que é necessário para o desenvolvimento gradual e pacífico da humanidade é que suas tendências não sejam perturbadas, nem tenham a liberdade de seus movimentos destruída.
\end{quote}

Desde o tempo de John Stuart Mill, economistas de autoridade têm estado em forte reação contra todas essas ideias. "Dificilmente um único economista inglês de renome", como o Professor Cannan expressou, "se juntará em um ataque frontal ao Socialismo em geral", embora, como ele também acrescenta, "quase todo economista, de renome ou não, esteja sempre pronto para apontar falhas na maioria das propostas socialistas".\footnote{\textit{Theories of Production and Distribution}, p. 494.} Os economistas não têm mais qualquer ligação com as filosofias teológicas ou políticas das quais o dogma da harmonia social nasceu, e sua análise científica não os leva a tais conclusões.

Cairnes, na palestra introdutória sobre "\textit{Political Economy and Laissez-faire}", que ele proferiu no University College, Londres, em 1870, foi talvez o primeiro economista ortodoxo a desferir um ataque frontal ao \textit{laissez-faire} em geral. "A máxima do \textit{laissez-faire}", declarou ele, "não tem base científica alguma, mas é na melhor das hipóteses uma mera regra prática útil".\footnote{Cairnes descreveu bem a "noção predominante" na seguinte passagem da mesma palestra: "A noção predominante é que a E.P. propõe-se a mostrar que a riqueza pode ser acumulada mais rapidamente e distribuída de forma mais justa; isto é, que o bem-estar humano pode ser promovido mais eficazmente, pelo simples processo de deixar as pessoas por si mesmas; deixando os indivíduos, ou seja, seguir os impulsos do auto-interesse, sem restrições do Estado ou da opinião pública, desde que se abstenham de força e fraude. Esta é a doutrina comumente conhecida como \textit{laissez-faire}; e consequentemente a economia política é, creio eu, muito geralmente considerada como uma espécie de renderização científica desta máxima — uma vindicação da liberdade de iniciativa individual e de contrato como a única e suficiente solução de todos os problemas industriais".} Isso, nos últimos cinquenta anos, tem sido a visão de todos os principais economistas. Parte do trabalho mais importante de Alfred Marshall — para citar um exemplo — foi dirigida à elucidação dos casos principais em que o interesse privado e o interesse social \textit{não} são harmoniosos. No entanto, a atitude cautelosa e não dogmática dos melhores economistas não prevaleceu contra a opinião geral de que um \textit{laissez-faire} individualista é tanto o que eles deveriam ensinar quanto o que, de fato, ensinam.

\begin{center}
\textbf{III}
\end{center}

Os economistas, como outros cientistas, escolheram a hipótese da qual partem, e que oferecem aos iniciantes, porque é a mais simples, e não porque seja a mais próxima dos fatos. Em parte por esta razão, mas em parte, admito, porque foram influenciados pelas tradições do assunto, eles começaram assumindo um estado de coisas em que a distribuição ideal dos recursos produtivos pode ser provocada por indivíduos agindo independentemente pelo método de tentativa e erro, de tal forma que aqueles indivíduos que se movem na direção certa destruirão pela concorrência aqueles que se movem na direção errada. Isso implica que não deve haver misericórdia ou proteção para aqueles que embarcam seu capital ou seu trabalho na direção errada. É um método de trazer os fazedores de lucro mais bem-sucedidos ao topo por uma luta implacável pela sobrevivência, que seleciona o mais eficiente pela falência do menos eficiente. Não conta o custo da luta, mas olha apenas para os benefícios do resultado final, que são assumidos como permanentes. Sendo o objetivo da vida colher as folhas dos galhos até a maior altura possível, a maneira mais provável de alcançar este fim é deixar as girafas com os pescoços mais longos matarem de fome aquelas cujos pescoços são mais curtos.

Correspondendo a este método de alcançar a distribuição ideal dos instrumentos de produção entre diferentes propósitos, há uma suposição semelhante quanto a como alcançar a distribuição ideal do que está disponível para consumo. Em primeiro lugar, cada indivíduo descobrirá o que, entre os possíveis objetos de consumo, ele mais deseja pelo método de tentativa e erro "na margem", e desta forma não apenas cada consumidor passará a distribuir seu consumo de forma mais vantajosa, mas cada objeto de consumo encontrará seu caminho para a boca do consumidor cujo gosto por ele é maior comparado com o dos outros, porque esse consumidor superará os demais nos lances. Assim, se apenas deixarmos as girafas por conta própria, (1) a quantidade máxima de folhas será colhida porque as girafas com os pescoços mais longos, à força de matar as outras de fome, chegarão mais perto das árvores; (2) cada girafa buscará as folhas que considerar mais suculentas entre as que estão ao seu alcance; e (3) as girafas cujo gosto por uma determinada folha for maior se esforçarão mais para alcançá-la. Desta forma, mais folhas e mais suculentas serão engolidas, e cada folha individual alcançará a garganta que julga merecer mais esforço.

Esta suposição, no entanto, de condições onde a seleção natural sem entraves leva ao progresso, é apenas uma das duas suposições provisórias que, tomadas como verdade literal, tornaram-se os pilares gêmeos do \textit{laissez-faire}. A outra é a eficácia, e de fato a necessidade, da oportunidade para a acumulação privada ilimitada de dinheiro como um incentivo ao esforço máximo. O lucro advém, sob o \textit{laissez-faire}, ao indivíduo que, seja por habilidade ou boa sorte, é encontrado com seus recursos produtivos no lugar certo e na hora certa. Um sistema que permite ao indivíduo habilidoso ou afortunado colher todos os frutos desta conjuntura evidentemente oferece um imenso incentivo à prática da arte de estar no lugar certo na hora certa. Assim, um dos motivos humanos mais poderosos, a saber, o amor ao dinheiro, é atrelado à tarefa de distribuir recursos econômicos da maneira melhor calculada para aumentar a riqueza.

O paralelismo entre o \textit{laissez-faire} econômico e o darwinismo, já brevemente observado, é agora visto, como Herbert Spencer foi o primeiro a reconhecer, como sendo de fato muito estreito. Assim como Darwin invocou o amor sexual, agindo através da seleção sexual, como um adjuvante da seleção natural pela competição, para direcionar a evolução ao longo de linhas que deveriam ser desejáveis bem como eficazes, o individualista invoca o amor ao dinheiro, agindo através da busca do lucro, como um adjuvante da seleção natural, para provocar a produção na maior escala possível do que é mais fortemente desejado conforme medido pelo valor de troca.

A beleza e a simplicidade de tal teoria são tão grandes que é fácil esquecer que ela não decorre dos fatos reais, mas de uma hipótese incompleta introduzida por uma questão de simplicidade. Além de outras objeções a serem mencionadas mais tarde, a conclusão de que indivíduos agindo independentemente para sua própria vantagem produzirão o maior agregado de riqueza depende de uma variedade de suposições irreais no sentido de que os processos de produção e consumo não são de forma alguma orgânicos, que existe um pré-conhecimento suficiente das condições e requisitos, e que existem oportunidades adequadas de obter este pré-conhecimento. Pois os economistas geralmente reservam para uma fase posterior de seu argumento as complicações que surgem — (1) quando as unidades eficientes de produção são grandes em relação às unidades de consumo, (2) quando custos fixos ou custos conjuntos estão presentes, (3) quando economias internas tendem à agregação da produção, (4) quando o tempo necessário para ajustes é longo, (5) quando a ignorância prevalece sobre o conhecimento, e (6) quando monopólios e combinações interferem na igualdade de negociação — eles reservam, ou seja, para uma fase posterior sua análise dos fatos reais. Além disso, muitos daqueles que reconhecem que a hipótese simplificada não corresponde exatamente ao fato concluem, no entanto, que ela representa o que é "natural" e, portanto, ideal. Eles consideram a hipótese simplificada como saúde e as complicações adicionais como doença.

No entanto, além desta questão de fato, existem outras considerações, familiares o suficiente, que justificadamente trazem para o cálculo o custo e o caráter da própria luta competitiva, e a tendência da riqueza ser distribuída onde ela não é mais apreciada. Se temos o bem-estar das girafas em mente, não devemos ignorar o sofrimento dos pescoços mais curtos que morrem de fome, ou as folhas doces que caem no chão e são pisoteadas na luta, ou o excesso de alimentação das girafas de pescoço longo, ou o olhar maligno de ansiedade ou ganância agressiva que anuvia as faces outrora suaves do rebanho.

Mas os princípios do \textit{laissez-faire} tiveram outros aliados além dos livros de economia. Deve-se admitir que eles foram confirmados nas mentes de pensadores sensatos e do público razoável pela má qualidade das propostas adversárias — o protecionismo de um lado e o socialismo marxista do outro. No entanto, essas doutrinas são ambas caracterizadas, não apenas ou principalmente por infringirem a presunção geral em favor do \textit{laissez-faire}, mas por mera falácia lógica. Ambas são exemplos de pensamento pobre, de incapacidade de analisar um processo e segui-lo até sua conclusão. Os argumentos contra elas, embora reforçados pelo princípio do \textit{laissez-faire}, não o exigem estritamente. Das duas, o protecionismo é pelo menos plausível, e as forças que contribuem para sua popularidade não são de se admirar. Mas o socialismo marxista deve sempre permanecer um prodígio para os historiadores da opinião — como uma doutrina tão ilógica e tão maçante pode ter exercido uma influência tão poderosa e duradoura sobre as mentes dos homens e, através delas, sobre os eventos da história. De qualquer forma, as óbvias deficiências científicas dessas duas escolas contribuíram grandemente para o prestígio e a autoridade do \textit{laissez-faire} do século XIX.

Nem a divergência mais notável para a ação social centralizada em grande escala — a condução da última guerra — encorajou reformadores ou dissipou preconceitos antiquados. Há muito a ser dito, é verdade, de ambos os lados. A experiência de guerra na organização da produção socializada deixou alguns observadores próximos otimisticamente ansiosos para repeti-la em condições de paz. O socialismo de guerra inquestionavelmente alcançou uma produção de riqueza em uma escala muito maior do que jamais conhecemos na paz, pois embora os bens e serviços entregues estivessem destinados à extinção imediata e infrutífera, não deixavam de ser riqueza. No entanto, a dissipação de esforço também foi prodigiosa, e a atmosfera de desperdício e de não contabilizar os custos era repugnante para qualquer espírito econômico ou previdente.

Finalmente, o individualismo e o \textit{laissez-faire} não poderiam, apesar de suas raízes profundas nas filosofias políticas e morais do final do século XVIII e início do XIX, ter assegurado seu domínio duradouro sobre a condução dos assuntos públicos se não fosse por sua conformidade com as necessidades e desejos do mundo dos negócios da época. Eles deram pleno escopo aos nossos heróis de outrora, os grandes homens de negócios. "Pelo menos metade da melhor habilidade no mundo ocidental", costumava dizer Marshall, "está engajada nos negócios". Uma grande parte da "imaginação superior" da época era assim empregada. Foi nas atividades desses homens que nossas esperanças de progresso se centraram.

\begin{quote}
Homens desta classe [escreveu Marshall]\footnote{"\textit{The Social Possibilities of Economic Chivalry}", \textit{Economic Journal}, XVII (1907), 9.} vivem em visões constantemente mutáveis, moldadas em seus próprios cérebros, de várias rotas para o fim desejado; das dificuldades que a Natureza lhes oporá em cada rota, e dos artifícios pelos quais esperam vencer a oposição dela. Esta imaginação ganha pouco crédito com o povo, porque não lhe é permitido correr solta; sua força é disciplinada por uma vontade mais forte; e sua glória suprema é ter alcançado grandes fins por meios tão simples que ninguém saberá, e ninguém além de especialistas sequer adivinhará, como uma dúzia de outros expedientes, cada um sugerindo tanto brilho ao observador apressado, foram postos de lado em favor dele. A imaginação de tal homem é empregada, como a do mestre jogador de xadrez, em prever os obstáculos que podem se opor ao sucesso de seus projetos de longo alcance, e constantemente rejeitando sugestões brilhantes porque ele imaginou para si mesmo os contra-ataques a elas. Sua força nervosa está no extremo oposto da natureza humana daquela irresponsabilidade nervosa que concebe esquemas utópicos apressados, e que deve ser comparada antes à facilidade audaciosa de um jogador fraco, que resolverá rapidamente o problema de xadrez mais difícil assumindo para si o direito de mover tanto as peças pretas quanto as brancas.
\end{quote}

Esta é uma bela imagem do grande capitão de indústria, o mestre-individualista, que nos serve servindo a si mesmo, exatamente como qualquer outro artista faz. No entanto, este, por sua vez, está se tornando um ídolo manchado. Ficamos mais duvidosos se é ele quem nos levará ao paraíso pela mão.

Esses muitos elementos contribuíram para o atual viés intelectual, a constituição mental, a ortodoxia do dia. A força convincente de muitas das razões originais desapareceu, mas, como de costume, a vitalidade das conclusões sobrevive a elas. Sugerir ação social para o bem público à \textit{City} de Londres é como discutir a \textit{Origem das Espécies} com um bispo sessenta anos atrás. A primeira reação não é intelectual, mas moral. Uma ortodoxia está em questão e, quanto mais persuasivos os argumentos, mais grave a ofensa. No entanto, aventurando-me no covil do monstro letárgico, pelo menos tracei suas reivindicações e linhagem de modo a mostrar que ele nos governou mais por direito hereditário do que por mérito pessoal.

\begin{center}
\textbf{§ IV}
\end{center}

Vamos limpar do terreno os princípios metafísicos ou gerais sobre os quais, de tempos em tempos, o \textit{laissez-faire} foi fundado. Não é verdade que os indivíduos possuem uma "liberdade natural" prescritiva em suas atividades econômicas. Não existe um "contrato" conferindo direitos perpétuos àqueles que Têm ou àqueles que Adquirem. O mundo não é governado do alto de forma que o interesse privado e o social sempre coincidam. Não é assim administrado aqui embaixo que na prática eles coincidam. Não é uma dedução correta dos princípios da economia que o auto-interesse esclarecido sempre opera no interesse público. Nem é verdade que o auto-interesse seja geralmente esclarecido; mais frequentemente os indivíduos que agem separadamente para promover seus próprios fins são ignorantes ou fracos demais para alcançar até mesmo estes. A experiência não mostra que os indivíduos, quando formam uma unidade social, são sempre menos clarividentes do que quando agem separadamente.

Não podemos, portanto, decidir em bases abstratas, mas devemos tratar em seus méritos detalhadamente o que Burke chamou de "um dos problemas mais delicados na legislação, a saber, determinar o que o Estado deve tomar para si para dirigir pela sabedoria pública, e o que deve deixar, com o mínimo de interferência possível, ao esforço individual".\footnote{Citado por McCulloch em seu \textit{Principles of Political Economy}.} Temos que discriminar entre o que Bentham, em sua nomenclatura esquecida mas útil, costumava denominar \textit{Agenda} e \textit{Não-Agenda}, e fazer isso sem a presunção prévia de Bentham de que a interferência é, ao mesmo tempo, "geralmente desnecessária" e "geralmente perniciosa".\footnote{\textit{Manual of Political Economy} de Bentham, publicado postumamente, na edição de Bowring (1843).} Talvez a tarefa principal dos economistas nesta hora seja distinguir novamente a \textit{Agenda} do governo da \textit{Não-Agenda}; e a tarefa companheira da política é idealizar formas de governo dentro de uma democracia que sejam capazes de realizar a \textit{Agenda}. Ilustrarei o que tenho em mente com dois exemplos.

(1) Acredito que, em muitos casos, o tamanho ideal para a unidade de controle e organização situa-se em algum lugar entre o indivíduo e o Estado moderno. Sugiro, portanto, que o progresso reside no crescimento e no reconhecimento de órgãos semi-autônomos dentro do Estado — órgãos cujo critério de ação dentro de seu próprio campo seja unicamente o bem público conforme o entendem, e de cujas deliberações os motivos de vantagem privada sejam excluídos, embora algum lugar ainda possa ser necessário deixar, até que o âmbito do altruísmo dos homens cresça mais, para a vantagem separada de grupos, classes ou faculdades particulares — corpos que no curso ordinário dos assuntos são principalmente autônomos dentro de suas limitações prescritas, mas estão sujeitos em última instância à soberania da democracia expressa através do Parlamento.

Proponho um retorno, pode-se dizer, em direção às concepções medievais de autonomias separadas. Mas, na Inglaterra, pelo menos, as corporações são um modo de governo que nunca deixou de ser importante e é simpático às nossas instituições. É fácil dar exemplos, do que já existe, de autonomias separadas que alcançaram ou estão se aproximando do modo que designo — as universidades, o Banco da Inglaterra, a Autoridade do Porto de Londres, até talvez as companhias ferroviárias. Na Alemanha, existem sem dúvida instâncias análogas.

Mas mais interessante do que estes é a tendência das instituições de capital aberto, quando atingem certa idade e tamanho, de se aproximarem do status de corporações públicas em vez do de empresa privada individualista. Um dos desenvolvimentos mais interessantes e despercebidos das décadas recentes tem sido a tendência da grande empresa de socializar-se. Chega um ponto no crescimento de uma grande instituição — particularmente uma grande ferrovia ou grande empresa de utilidade pública, mas também um grande banco ou uma grande companhia de seguros — em que os proprietários do capital, ou seja, os acionistas, estão quase inteiramente dissociados da gestão, com o resultado de que o interesse pessoal direto destes últimos na obtenção de grandes lucros torna-se secundário. Quando este estágio é alcançado, a estabilidade geral e a reputação da instituição são mais consideradas pela gestão do que o máximo de lucro para os acionistas. Os acionistas devem ser satisfeitos por dividendos convencionalmente adequados; mas, uma vez que isso é garantido, o interesse direto da gestão frequentemente consiste em evitar críticas do público e dos clientes da empresa. Este é particularmente o caso se o seu grande tamanho ou posição semi-monopolística os torna conspícuos aos olhos do público e vulneráveis ao ataque público. O exemplo extremo, talvez, desta tendência no caso de uma instituição, teoricamente a propriedade irrestrita de pessoas privadas, é o Banco da Inglaterra. É quase verdade dizer que não há classe de pessoas no reino de quem o Governador do Banco da Inglaterra pense menos quando decide sobre sua política do que de seus acionistas. Seus direitos, além de seu dividendo convencional, já afundaram para as vizinhanças de zero. Mas a mesma coisa é parcialmente verdadeira de muitas outras grandes instituições. Elas estão, com o passar do tempo, socializando-se.

Não que isso seja ganho puro. As mesmas causas promovem o conservadorismo e um declínio da iniciativa. De fato, já temos nesses casos muitas das falhas, bem como as vantagens do Socialismo de Estado. No entanto, vemos aqui, creio eu, uma linha natural de evolução. A batalha do Socialismo contra o lucro privado ilimitado está sendo vencida em detalhes, hora após hora. Nestes campos específicos — permanece aguda em outros lugares — este não é mais o problema premente. Não há, por exemplo, nenhuma questão política dita importante tão realmente sem importância, tão irrelevante para a reorganização da vida econômica da Grã-Bretanha, quanto a nacionalização das ferrovias.

É verdade que muitos grandes empreendimentos, particularmente empresas de utilidade pública e outros negócios que exigem um grande capital fixo, ainda precisam ser semi-socializados. Mas devemos manter nossas mentes flexíveis quanto às formas desse semi-socialismo. Devemos tirar total vantagem das tendências naturais do dia, e provavelmente devemos preferir corporações semi-autônomas a órgãos do governo central pelos quais os ministros de Estado são diretamente responsáveis.

Critico o Socialismo de Estado doutrinário, não porque ele busque engajar os impulsos altruístas dos homens a serviço da sociedade, ou porque ele se afaste do \textit{laissez-faire}, ou porque ele retire do homem a sua liberdade natural de ganhar um milhão, ou porque ele tenha coragem para experimentos ousados. Todas essas coisas eu aplaudo. Critico-o porque ele perde o significado do que está realmente acontecendo; porque ele é, de fato, pouco melhor que uma sobrevivência empoeirada de um plano para enfrentar os problemas de cinquenta anos atrás, baseado em um mal-entendido do que alguém disse cem anos atrás. O Socialismo de Estado do século XIX surgiu de Bentham, da livre concorrência, etc., e é em alguns aspectos uma versão mais clara, em outros uma versão mais confusa, da mesma filosofia que fundamenta o individualismo do século XIX. Ambos igualmente colocaram todo o seu estresse na liberdade, um negativamente para evitar limitações à liberdade existente, o outro positivamente para destruir monopólios naturais ou adquiridos. São reações diferentes à mesma atmosfera intelectual.

(2) Venho agora a um critério de \textit{Agenda} que é particularmente relevante para o que é urgente e desejável fazer num futuro próximo. Devemos visar a separar aqueles serviços que são \textit{tecnicamente sociais} daqueles que são \textit{tecnicamente individuais}. A \textit{Agenda} mais importante do Estado refere-se não àquelas atividades que os indivíduos privados já estão cumprindo, mas àquelas funções que caem fora da esfera do indivíduo, àquelas decisões que são tomadas por \textit{ninguém} se o Estado não as tomar. O importante para o governo não é fazer coisas que os indivíduos já estão fazendo, e fazê-las um pouco melhor ou um pouco pior; mas fazer aquelas coisas que atualmente não são feitas de forma alguma.

Não cabe ao meu propósito nesta ocasião desenvolver políticas práticas. Limito-me, portanto, a citar alguns exemplos do que quero dizer a partir daqueles problemas sobre os quais calhou de eu ter pensado mais.

Muitos dos maiores males econômicos do nosso tempo são frutos do risco, da incerteza e da ignorância. É porque indivíduos particulares, afortunados em situação ou habilidades, são capazes de tirar proveito da incerteza e da ignorância, e também porque, pela mesma razão, os grandes negócios são frequentemente uma loteria, que grandes desigualdades de riqueza surgem; e esses mesmos fatores são também a causa do desemprego do trabalho, ou do desapontamento de expectativas comerciais razoáveis, e do comprometimento da eficiência e da produção. No entanto, a cura reside fora das operações dos indivíduos; pode até ser do interesse dos indivíduos agravar a doença. Acredito que a cura para estas coisas deve ser buscada em parte no controle deliberado da moeda e do crédito por uma instituição central, e em parte na coleta e disseminação em grande escala de dados relativos à situação comercial, incluindo a plena publicidade, por lei se necessário, de todos os fatos comerciais que seja útil conhecer. Estas medidas envolveriam a sociedade no exercício da inteligência diretiva através de algum órgão apropriado de ação sobre muitas das intrincadas engrenagens íntimas dos negócios privados, mas deixariam a iniciativa e o empreendimento privados desimpedidos. Mesmo que estas medidas se mostrem insuficientes, elas nos fornecerão um melhor conhecimento do que temos agora para dar o próximo passo.

Meu segundo exemplo refere-se à poupança e ao investimento. Acredito que algum ato coordenado de julgamento inteligente é necessário quanto à escala em que é desejável que a comunidade como um todo deva poupar, a escala em que estas poupanças devam ir para o estrangeiro sob a forma de investimentos externos, e se a atual organização do mercado de capitais distribui as poupanças ao longo dos canais mais produtivos nacionalmente. Não creio que estes assuntos devam ser deixados inteiramente ao sabor do julgamento privado e dos lucros privados, como são atualmente.

Meu terceiro exemplo diz respeito à população. Já chegou o tempo em que cada país necessita de uma política nacional ponderada sobre qual o tamanho da população, seja maior ou menor que a atual ou a mesma, que é mais conveniente. E tendo estabelecido esta política, devemos dar passos para colocá-la em operação. Poderá chegar o tempo, um pouco mais tarde, em que a comunidade como um todo deva prestar atenção à qualidade inata, bem como aos meros números, de seus futuros membros.

\begin{center}
\textbf{§ V}\footnote{O número do capítulo não apareceu, é claro, na edição original de \textit{Essays in Persuasion} — Ed.}
\end{center}

Estas reflexões foram dirigidas para possíveis melhorias na técnica do capitalismo moderno pela agência da ação coletiva. Não há nada nelas que seja seriamente incompatível com o que me parece ser a característica essencial do capitalismo, a saber, a dependência de um apelo intenso aos instintos de ganhar dinheiro e de amar o dinheiro dos indivíduos como a principal força motriz da máquina econômica. Nem devo eu, tão perto do meu fim, desviar-me para outros campos. No entanto, farei bem em lembrar-lhes, em conclusão, que as disputas mais ferozes e as divisões de opinião mais profundamente sentidas provavelmente serão travadas nos próximos anos não em torno de questões técnicas, onde os argumentos de ambos os lados são principalmente econômicos, mas em torno daquelas que, por falta de palavras melhores, podem ser chamadas de psicológicas ou, talvez, morais.

Na Europa, ou pelo menos em algumas partes da Europa — mas não, creio eu, nos Estados Unidos da América — existe uma reação latente, um tanto generalizada, contra basear a sociedade na medida em que o fazemos no fomento, encorajamento e proteção dos motivos pecuniários dos indivíduos. Uma preferência por organizar nossos assuntos de tal forma que se apele ao motivo do lucro o menos possível, em vez de o máximo possível, não precisa ser inteiramente \textit{a priori}, mas pode ser baseada na comparação de experiências. Diferentes pessoas, de acordo com a sua escolha de profissão, descobrem que o motivo do lucro desempenha um papel grande ou pequeno em suas vidas diárias, e os historiadores podem nos contar sobre outras fases da organização social em que este motivo desempenhou um papel muito menor do que desempenha agora. A maioria das religiões e a maioria das filosofias deprecam, para dizer o mínimo, um modo de vida influenciado principalmente por considerações de lucro pessoal em dinheiro. Por outro lado, a maioria dos homens hoje rejeita noções ascéticas e não duvida das vantagens reais da riqueza. Além disso, parece-lhes óbvio que não se pode passar sem o motivo do lucro e que, exceto por certos abusos admitidos, ele faz bem o seu trabalho. Em resultado, o homem comum desvia a sua atenção do problema e não tem uma ideia clara do que realmente pensa e sente sobre toda esta maldita questão.

A confusão de pensamento e de sentimento leva à confusão de fala. Muitas pessoas, que estão realmente objetando ao capitalismo como um modo de vida, argumentam como se estivessem objetando a ele com base na sua ineficiência em atingir os seus próprios objetivos. Pelo contrário, os devotos do capitalismo são frequentemente indevidamente conservadores, e rejeitam reformas na sua técnica, que poderiam realmente fortalecê-lo e preservá-lo, por medo de que possam provar ser os primeiros passos para longe do próprio capitalismo. No entanto, um tempo pode estar chegando em que ficaremos mais esclarecidos do que atualmente sobre quando estamos falando do capitalismo como uma técnica eficiente ou ineficiente, e quando estamos falando dele como desejável ou objetável em si mesmo. Pela minha parte, penso que o capitalismo, sabiamente gerido, pode provavelmente ser tornado mais eficiente para atingir fins econômicos do que qualquer sistema alternativo ainda à vista, mas que em si mesmo é em muitos aspectos extremamente objetável. O nosso problema é elaborar uma organização social que seja tão eficiente quanto possível sem ofender as nossas noções de um modo de vida satisfatório.

O próximo passo em frente deve vir, não da agitação política ou de experiências prematuras, mas do pensamento. Precisamos, por um esforço da mente, elucidar os nossos próprios sentimentos. Atualmente, a nossa simpatia e o nosso julgamento estão propensos a estar em lados diferentes, o que é um estado de espírito doloroso e paralisante. No campo da ação, os reformadores não serão bem-sucedidos até que possam perseguir firmemente um objeto claro e definido com os seus intelectos e os seus sentimentos em sintonia. Não há nenhum partido no mundo atualmente que me pareça estar a perseguir objetivos certos por métodos certos. A pobreza material fornece o incentivo para mudar precisamente em situações onde há muito pouca margem para experiências. A prosperidade material remove o incentivo exatamente quando poderia ser seguro arriscar-se. A Europa carece de meios, a América de vontade, para fazer um movimento. Precisamos de um novo conjunto de convicções que surjam naturalmente de um exame cândido dos nossos próprios sentimentos interiores em relação aos fatos exteriores.

\end{document}